\documentclass[instructions.tex]{subfiles}
\begin{document}
 
\chapter{Du vice contraire à la chasteté.}
 
\primo Si une princesse richement parée allait se jeter dans un cloaque, pour y caresser un pourceau et jouer avec ce vil animal, ce serait une conduite bien extravagante. O âme chrétienne ! vous êtes la princesse de l’univers, parée des ornemens de la grâce, quel est votre égarement de vous plonger dans le bourbier des sales voluptés, d’être le jouet de votre corps, qui est comme un animal !
 
Que dirait-on d’une reine qui obéirait à sa servante, jusqu’à se laisser mettre sous ses pieds ? Votre conduite est encore plus honteuse. Votre corps doit vous être soumis, et vous devenez l’esclave de ses fantaisies ; ce corps est un vil esclave, et vous en faites votre maître ; ce corps est votre ennemi, et vous le flattez ; mais bientôt il sera votre tyran, et ses plaisirs vous coûteront cher.
 
Si un homme avait la témérité d’entrer dans le temple de Dieu, pour y défigurer les images du Sauveur; s’il plaçait dans le sanctuaire, où repose Jésus-Christ, la figure d’un animal, ou une idole, pour l’y adorer; s’il profanait indignement les vases sacrés, y aurait-il un supplice assez rigoureux pour le punir ? Ne commettez-vous pas, en quelque façon, un pareil attentat? Vous êtes, dit saint Paul, les temples de Dieu qui habite en vous. Vos membres sont les membres de Jésus-Christ; et vous n’avez point d’horreur de souiller dans vous ou dans les autres ces membres sacrés ! Votre âme est l’image de Dieu, elle est teinte du sang du Sauveur, et vous la défigurez par ce vice honteux ! Le Saint-Esprit est dans votre cœur comme dans son temple; vous l’en chassez pour y placer une idole de chair, pour y faire régner le plus infâme des maîtres, qui est le démon de l’impureté.
 
\secundo Oh! que vous êtes aveugle, si vous ne comprenez pas combien ce vice est détestable ! Voyez les châtimens dont Dieu l’a puni. N’a-t-il pas fait périr dans les eaux de sa colère tous les hommes hors de l’arche, pour leurs impuretés; fait pleuvoir le feu et le soufre pour consumer les impudiques habitans de Sodome; puni de mort le misérable Onan, qui profanait avec son épouse la sainteté du mariage; fait massacrer plus de vingt-trois mille Israélites en punition de leurs dissolutions?
 
Écoutez saint Paul, qui dit que Dieu abandonne les impudiques à leur sens réprouvé, et que le sort ordinaire de ces infâmes est le désespoir\footnote{Desperantes semetipsos tradiderunt impudicitiae. Ephe.4.19.}. Écoutez ce grand apôtre, qui vous dit de la part de Jésus-Christ, que les fornicateurs et les impudiques n’entreront jamais dans le royaume des cieux. Croyez-en ce que vous voudrez; mais sachez que si vous ne pleurez vos impudicités, si vous ne quittez vos habitudes et vos criminelles amitiés, vous les pleurerez un jour dans les feux éternels. Les gens de bien et les personnes chastes ont horreur de votre conduite, et ne peuvent comprendre l’aveuglement qui vous transporte. Bien loin de gémir et d’avoir horreur de vous-même, vous ne faites que rire et plaisanter de vos infamies; semblable à ces deux voluptueux dont parle l’Écriture, qui se renversèrent l’esprit pour étouffer en eux tout sentiment de crainte de Dieu \footnote{Everterunt sensum suum. Dan. 13.}.
 
En effet, l’impudique est dans un endurcissement si profond, qu’il étouffe tout sentiment de pudeur, que rien ne le touche ne l’arrête. Ses regards sont autant d’adultères qu’il commet dans son cœur \footnote{Oculos Habentes plenos adulterii.}. Ses discours sont autant d’étincelles qui jettent partout le feu impur. Ses pensées, ses désirs, ne se portent qu’aux objets de sa passion. Dans le lieu saint, où les démons mêmes n’entrent qu’en tremblant, il porte son idole dans son cœur, et l’adore au préjudice du vrai Dieu, qu’il fait semblant d’adorer.
 
L’ignominie et l’horreur de sa conduite, l’infamie qu’il attire sur sa famille, l’opprobre et la désolation des vierges qu’il a flétries et qu’il a perdues, les redoutables jugemens de Dieu qui l’atteindent, tout cela ne fait point d’impression sur son cœur abruti. Il pèche seul; il pèche en compagnie; il pèche sans remords; il pèche partout; il pèche sans cesse, dit le Saint-Esprit\footnote{Incessabilis delicti 2. Petr. 2.}. O mon Dieu! que l’homme est misérable, lorsqu’il suit les désirs de sa chair, et qu’il ne sent plus l’horreur et la honte de ce vice !
 
\tertio Malheureux, pourquoi perdez-vous votre âme pour les infâmes plaisirs de votre corps ? comprenez-vous quelle est la vilité de ce corps, et quelle est la dignité de votre âme ? Votre corps n’est qu’une vile partie de vous-même, un amas de terre, et comme un habit de boue que vous portez; mais votre âme c’est vous-même, et la plus noble partie de vous-même. Votre corps tombera en pourriture, votre âme vivra toujours. Quelle fureur de perdre cette âme précieuse, pour un corps qui n’est que boue et que poussière.
 
Considérez un corps lorsque l’âme qui en faisait toute la force et la beauté en est séparée, que verrez-vous ? un cadavre hideux, infect, horrible. Celui de la personne qui vous charme, aussi bien que le vôtre, sera bientôt de même. C’est cependant ce vil cadavre que vous flattez, pour lequel vous perdez votre âme et votre Dieu. O aveuglement!
 
Qu’est-ce que la beauté de ce jeune homme, de cette fille qui vous enchante ? Que sont les appas de ce visage qui vous plaît, qui est l’idole de vos yeux, que vous caressez par tant de libertés folâtres? c’est une fleur qui sera bientôt flétrie. Il ne faut qu’une contusion, un ulcère, une défaillance de cœur, pour changer ce visage, et le rendre difforme à faire horreur. O folie des enfans des hommes, d’attacher son cœur à si peu de chose : et de perdre Dieu pour des objets si fragiles!
 
Prenez en main une tête de mort; voyez cette tête affreuse, sans cheveux, les yeux et les oreilles creusés, les joues et les dents décharnées; voilà l’état où sera bientôt ce visage que vous regardez avec tant de complaisance. Un jeune homme, inconsolable de la mort d’une fille, ne pouvait en perdre le souvenir. Un de ses amis l’ayant conduit où elle était enterrée, fit lever la pierre du tombeau et les suaires qui la couvraient. Dans ce moment une insupportable puanteur faillit les suffoquer. Les vers et la pourriture sortaient déjà de la bouche et des yeux de ce cadavre. Ce jeune homme jeta un cri d’horreur, et voulut s’enfuir. De quoi avez-vous peur? lui dit son ami; approchez-vous. Voilà le visage de cette fille que vous avez tant aimée et caressée, qui pleure à présent en l’autre monde les péchés que vous lui avez fait commettre; apprenez, à la vue de ces objets, à ne plus attacher votre cœur à des choses si indignes de vous. Ce jeune homme profita de cet avis et se convertit.
 
\section{Résolutions}
 
\primo Je regarderai désormais le vice opposé à la chasteté comme quelque sorte sacrilége du temple de Dieu et des membres de Jésus-Christ, comme la source féconde de péchés, de désordres et de crimes, qui conduisent à la damnation éternelle.
 
\secundo Pour en concevoir de plus en plus de l’éloignement, je me rappellerai souvent les châtimens terribles qu’il a fait tomber sur tant d’hommes qui s’en laissaient dominer, sur des villes entières, sur des royaumes florissans, enfin sur toute la terre.

\prayer{O mon Dieu! je reconnois humblement ma foiblesse : ce n’est pas de mes propres forces, mais de votre grâce, que j’attends la victoire dans les combats que j’aurai désormais à soutenir contre les convoitises de la chair; aidez-moi, soutenez-moi, sauvez-moi.}
 
\end{document}
